\documentclass[12pt]{article}
\usepackage{amsmath}
\usepackage{amssymb}
\usepackage{geometry}
\usepackage{graphicx}
\usepackage{cancel}
\geometry{letterpaper, margin=1in}

\title{\textbf{Homework  1.2: Solutions and Initial Value Problems}}
\author{Tashfeen Omran}
\date{September 2026}

\begin{document}

\maketitle
\tableofcontents
\newpage

\section{Part 1: Comprehensive Introduction, Context, and Prerequisites}

\subsection{Core Concepts}
A \textbf{Differential Equation (DE)} is an equation that relates an unknown function to one or more of its derivatives. While standard algebraic equations ask you to solve for a specific variable (a number), differential equations ask you to solve for a \textit{function}. 
\begin{itemize}
    \item \textbf{Ordinary Differential Equation (ODE):} An equation involving an unknown function of a single variable and its ordinary derivatives (e.g., $\frac{dy}{dx} = y + x$).
    \item \textbf{Order:} The highest derivative present in the equation. For example, $y'' + y = 0$ is a second-order ODE.
    \item \textbf{Solution to a DE:} A function $\phi(x)$ that, when substituted into the DE, yields a true identity on some interval $I$. 
    \item \textbf{Explicit vs. Implicit Solutions:} An \textit{explicit} solution gives the dependent variable isolated on one side (e.g., $y = x^2 + C$). An \textit{implicit} solution defines the relationship between $x$ and $y$ without isolating $y$ (e.g., $x^2 + y^2 = C$).
    \item \textbf{Initial Value Problem (IVP):} A differential equation paired with one or more initial conditions (e.g., $y(x_0) = y_0$). These conditions allow us to find a specific solution (a particular solution) from the general family of solutions.
\end{itemize}

\subsection{Intuition and Derivation}
In the real world, it is often much easier to measure the \textit{rate of change} of a system than to measure the system's exact state at any given time. For instance, we know from Newton's Law of Cooling that the rate at which an object cools ($\frac{dT}{dt}$) is proportional to the difference between the object's temperature ($T$) and the ambient temperature ($T_a$). This is naturally written as $\frac{dT}{dt} = -k(T - T_a)$. The differential equation arises organically from the physical observation. Solving the equation essentially means "undoing" the rate of change to find the original state function $T(t)$, which requires continuous summation (integration).

\subsection{Historical Context and Motivation}
The foundations of differential equations were laid in the 17th century by Isaac Newton and Gottfried Wilhelm Leibniz concurrently with their invention of calculus. Newton originally called them "fluxional equations." They were primarily motivated by physics and astronomy: predicting planetary orbits, modeling the motion of pendulums, and understanding classical mechanics. Later, mathematicians like Euler, Bernoulli, and Laplace expanded these techniques to solve increasingly complex problems in fluid dynamics, heat transfer, and electromagnetism.

\subsection{Key Formulas and Theorems}
\textbf{1. The Existence and Uniqueness Theorem (for first-order IVPs):}\\
Consider the initial value problem:
\[ \frac{dy}{dx} = f(x, y), \quad y(x_0) = y_0 \]
If $f(x, y)$ and its partial derivative with respect to $y$, $\frac{\partial f}{\partial y}$, are both continuous on some rectangular region $R$ in the $xy$-plane that contains the point $(x_0, y_0)$, then there exists a \textit{unique} solution $y(x)$ to the IVP on some interval $I$ containing $x_0$.

\subsection{Prerequisites}
To master this topic, you must be proficient in:
\begin{itemize}
    \item \textbf{Derivative Rules:} Power rule, product rule, quotient rule, and especially the \textbf{chain rule}.
    \item \textbf{Implicit Differentiation:} Taking the derivative of both sides of an equation with respect to $x$ and solving for $\frac{dy}{dx}$.
    \item \textbf{Partial Derivatives:} Treating all variables except the one of interest as constants.
    \item \textbf{Algebraic Manipulations:} Factoring, simplifying complex fractions, and working with exponents and radicals.
\end{itemize}

\newpage
\section{Part 2: Detailed Homework Solutions}

\subsection*{Problem 1}
\textbf{Part (a):} Show that $y^8 + x - 5 = 0$ is an implicit solution to $\frac{dy}{dx} = -\frac{1}{8y^7}$ on the interval $(-\infty, 5)$. Convert the implicit solution to explicit functions and find their derivatives. Explain the interval constraint.

\textbf{Solution:}
First, verify the implicit solution by differentiating both sides of $y^8 + x - 5 = 0$ with respect to $x$:
\begin{align*}
    \frac{d}{dx} \left( y^8 + x - 5 \right) &= \frac{d}{dx} (0) \\
    8y^7 \cdot \frac{dy}{dx} + 1 - 0 &= 0 \\
    8y^7 \frac{dy}{dx} &= -1 \\
    \frac{dy}{dx} &= -\frac{1}{8y^7}
\end{align*}
This matches the given differential equation. Next, convert to explicit functions by solving for $y$:
\begin{align*}
    y^8 &= 5 - x \\
    y &= \pm \sqrt[8]{5 - x}
\end{align*}
Thus, the explicit solutions are $y = \sqrt[8]{5 - x}$ and $y = -\sqrt[8]{5 - x}$. 
Now, find their derivatives. Note that $\sqrt[8]{5-x} = (5-x)^{1/8}$:
\begin{align*}
    \frac{dy}{dx} &= \frac{d}{dx} \left[ (5-x)^{1/8} \right] = \frac{1}{8}(5-x)^{-7/8} \cdot (-1) = -\frac{1}{8(5-x)^{7/8}} = \frac{-1}{8\sqrt[8]{(5-x)^7}} \\
    \frac{dy}{dx} &= \frac{d}{dx} \left[ -(5-x)^{1/8} \right] = -\frac{1}{8}(5-x)^{-7/8} \cdot (-1) = \frac{1}{8\sqrt[8]{(5-x)^7}}
\end{align*}
\textit{Check Interval Validity:} At $x = 5$, $y = 0$, making $\frac{dy}{dx}$ undefined (division by zero). For $x > 5$, $5-x < 0$, and the 8th root of a negative number is not a real number. Therefore, neither $y$ nor $\frac{dy}{dx}$ are real-valued for $x > 5$, and $\frac{dy}{dx}$ is undefined at $x = 5$.
\textbf{Correct Choice:} A

\vspace{0.5cm}

\textbf{Part (b):} Show that $xy^7 - xy^7\cos x = 1$ is an implicit solution to $\frac{dy}{dx} = \frac{(-x\sin x + \cos x - 1)y}{7(x - x\cos x)}$. Convert to an explicit function, find its derivative, and explain interval constraints.

\textbf{Solution:}
First, factor out $xy^7$ to simplify:
\[ xy^7(1 - \cos x) = 1 \]
Differentiate both sides implicitly with respect to $x$ using the product rule. Let $u = x(1-\cos x)$ and $v = y^7$:
\begin{align*}
    \frac{d}{dx} \left[ x(1-\cos x)y^7 \right] &= \frac{d}{dx}(1) \\
    \left[ 1(1-\cos x) + x(\sin x) \right]y^7 + x(1-\cos x) \left[ 7y^6 \frac{dy}{dx} \right] &= 0
\end{align*}
Isolate $\frac{dy}{dx}$:
\begin{align*}
    7xy^6(1-\cos x) \frac{dy}{dx} &= -y^7 (1 - \cos x + x\sin x) \\
    \frac{dy}{dx} &= \frac{-y^7 (1 - \cos x + x\sin x)}{7xy^6(1-\cos x)} \\
    \frac{dy}{dx} &= \frac{-y (1 - \cos x + x\sin x)}{7x(1-\cos x)} = \frac{(-x\sin x + \cos x - 1)y}{7(x - x\cos x)}
\end{align*}
This matches the given differential equation. Next, convert to an explicit function:
\begin{align*}
    xy^7(1 - \cos x) &= 1 \\
    y^7 &= \frac{1}{x(1 - \cos x)} \\
    y &= \frac{1}{\sqrt[7]{x(1 - \cos x)}}
\end{align*}
To find the derivative of the explicit solution, write $y = (x - x\cos x)^{-1/7}$ and apply the chain rule:
\begin{align*}
    \frac{dy}{dx} &= -\frac{1}{7}(x - x\cos x)^{-8/7} \cdot \frac{d}{dx}(x - x\cos x) \\
    &= -\frac{1}{7}(x(1-\cos x))^{-8/7} \cdot (1 - (\cos x - x\sin x)) \\
    &= \frac{-(1 - \cos x + x\sin x)}{7 (x(1-\cos x))^{8/7}} = \frac{1 - \cos x + x\sin x}{-7\sqrt[7]{(x(1-\cos x))^8}}
\end{align*}
\textit{Check Interval Validity:} The denominators for both $y$ and $\frac{dy}{dx}$ contain the term $x(1-\cos x)$. This term equals $0$ when $x = 0$ and when $\cos x = 1$ (e.g., $x = 2\pi$). Division by zero makes both functions undefined.
\textbf{Correct Choices:} A ($\frac{dy}{dx}$ is undefined at $x=0$ and $x=2\pi$) and C ($y$ is undefined at $x=0$ and $x=2\pi$).

\subsection*{Problem 2}
\textbf{Determine whether the given function is a solution:} $y = \cos x + x^7$ for $\frac{d^2y}{dx^2} + y = x^7 + 42x^5$.

\textbf{Solution:}
Find the first and second derivatives of the given function $y$:
\begin{align*}
    y &= \cos x + x^7 \\
    \frac{dy}{dx} &= -\sin x + 7x^6 \\
    \frac{d^2y}{dx^2} &= -\cos x + 42x^5
\end{align*}
Substitute $y$ and $\frac{d^2y}{dx^2}$ into the left side of the differential equation:
\begin{align*}
    \text{LHS} &= \frac{d^2y}{dx^2} + y \\
    &= (-\cos x + 42x^5) + (\cos x + x^7) \\
    &= x^7 + 42x^5
\end{align*}
Since the LHS equals the RHS ($x^7 + 42x^5$), the function is indeed a solution.
\textbf{Answers to fill in:} $\frac{d^2y}{dx^2} = -\cos x + 42x^5$.

\subsection*{Problem 3}
\textbf{Determine whether the given function is a solution:} $\theta = 4e^{4t} - 2e^{5t}$ for $\frac{d^2\theta}{dt^2} - \theta + 30 = -13e^{5t}$.

\textbf{Solution:}
Find the first and second derivatives of $\theta$:
\begin{align*}
    \theta &= 4e^{4t} - 2e^{5t} \\
    \frac{d\theta}{dt} &= 16e^{4t} - 10e^{5t} \\
    \frac{d^2\theta}{dt^2} &= 64e^{4t} - 50e^{5t}
\end{align*}
Substitute $\theta$ and $\frac{d^2\theta}{dt^2}$ into the left side of the differential equation:
\begin{align*}
    \text{LHS} &= \frac{d^2\theta}{dt^2} - \theta + 30 \\
    &= (64e^{4t} - 50e^{5t}) - (4e^{4t} - 2e^{5t}) + 30 \\
    &= 60e^{4t} - 48e^{5t} + 30
\end{align*}
This LHS ($60e^{4t} - 48e^{5t} + 30$) does \textit{not} equal the RHS ($-13e^{5t}$). Therefore, it is \textbf{not} a solution.
\textbf{Answers to fill in:} $\frac{d\theta}{dt} = 16e^{4t} - 10e^{5t}$ and $\frac{d^2\theta}{dt^2} = 64e^{4t} - 50e^{5t}$.

\subsection*{Problem 4}
\textbf{Verify implicit solution via implicit differentiation:} $e^{6xy} + 2y^3 = 2x^3 + 1$ for $\frac{dy}{dx} = \frac{x^2e^{-6xy} - y}{y^2e^{-6xy} + x}$.

\textbf{Solution:}
Differentiate both sides with respect to $x$. Use the chain rule and product rule for $e^{6xy}$:
\begin{align*}
    \frac{d}{dx} \left[ e^{6xy} \right] + \frac{d}{dx} \left[ 2y^3 \right] &= \frac{d}{dx} \left[ 2x^3 + 1 \right] \\
    e^{6xy} \left( 6y + 6x \frac{dy}{dx} \right) + 6y^2 \frac{dy}{dx} &= 6x^2
\end{align*}
Divide the entire equation by 6 to simplify:
\begin{align*}
    e^{6xy} \left( y + x \frac{dy}{dx} \right) + y^2 \frac{dy}{dx} &= x^2 \\
    ye^{6xy} + xe^{6xy}\frac{dy}{dx} + y^2\frac{dy}{dx} &= x^2
\end{align*}
Isolate $\frac{dy}{dx}$:
\begin{align*}
    \frac{dy}{dx} \left( xe^{6xy} + y^2 \right) &= x^2 - ye^{6xy} \\
    \frac{dy}{dx} &= \frac{x^2 - ye^{6xy}}{xe^{6xy} + y^2}
\end{align*}
To match the given form, multiply the numerator and the denominator by $e^{-6xy}$:
\begin{align*}
    \frac{dy}{dx} &= \frac{(x^2 - ye^{6xy})e^{-6xy}}{(xe^{6xy} + y^2)e^{-6xy}} \\
    \frac{dy}{dx} &= \frac{x^2e^{-6xy} - y}{x + y^2e^{-6xy}}
\end{align*}
This verifies the solution.

\subsection*{Problem 5}
\textbf{Existence and Uniqueness Theorem Application:} $\frac{dy}{dx} = y^4 + x^2, \quad y(0)=2$.

\textbf{Solution:}
Let $f(x, y) = y^4 + x^2$. We must check the continuity of $f(x,y)$ and its partial derivative with respect to $y$:
\begin{align*}
    \frac{\partial f}{\partial y} &= \frac{\partial}{\partial y} (y^4 + x^2) = 4y^3
\end{align*}
Both $f(x, y) = y^4 + x^2$ and $\frac{\partial f}{\partial y} = 4y^3$ are polynomials in two variables. Polynomials are continuous everywhere (on the entire $xy$-plane). Because both functions are continuous in any rectangle containing the initial point $(0, 2)$, the theorem \textit{implies} the existence of a unique solution.
\textbf{Correct Choice:} B. Both are continuous.

\subsection*{Problem 6}
\textbf{Existence and Uniqueness Theorem Application:} $y^4 \frac{dy}{dx} = x^5, \quad y(1)=0$.

\textbf{Solution:}
First, rewrite the equation in standard form $\frac{dy}{dx} = f(x,y)$:
\[ \frac{dy}{dx} = \frac{x^5}{y^4} \]
Here, $f(x,y) = \frac{x^5}{y^4}$. Notice that $f(x,y)$ is undefined, and thus discontinuous, when $y = 0$. 
The initial condition is $y(1) = 0$, meaning the point of interest is $(x_0, y_0) = (1, 0)$. Because the function $f(x,y)$ is discontinuous at $y=0$, we cannot draw any rectangle around $(1,0)$ where the function remains continuous. Thus, the theorem does \textit{not} imply the existence of a unique solution.
\textbf{Correct Choice:} C. $f(x,y)$ is not continuous.

\subsection*{Problem 7}
\textbf{Verify family of solutions and graph:} $\phi(x) = \frac{4}{1-ce^x}$ for $\frac{dy}{dx} = \frac{y^2 - 4y}{4}$.

\textbf{Solution:}
First, find $\phi'(x)$ using the chain rule. Let $\phi(x) = 4(1-ce^x)^{-1}$:
\begin{align*}
    \frac{dy}{dx} &= -4(1-ce^x)^{-2} \cdot (-ce^x) = \frac{4ce^x}{(1-ce^x)^2}
\end{align*}
Next, substitute $\phi(x)$ into the right side of the differential equation:
\begin{align*}
    \text{RHS} &= \frac{y^2 - 4y}{4} = \frac{\left( \frac{4}{1-ce^x} \right)^2 - 4\left( \frac{4}{1-ce^x} \right)}{4} \\
    &= \frac{1}{4} \left( \frac{16}{(1-ce^x)^2} - \frac{16}{1-ce^x} \right) \\
    &= \frac{4}{(1-ce^x)^2} - \frac{4}{1-ce^x}
\end{align*}
Find a common denominator of $(1-ce^x)^2$:
\begin{align*}
    \text{RHS} &= \frac{4}{(1-ce^x)^2} - \frac{4(1-ce^x)}{(1-ce^x)^2} \\
    &= \frac{4 - 4 + 4ce^x}{(1-ce^x)^2} \\
    &= \frac{4ce^x}{(1-ce^x)^2}
\end{align*}
Since LHS = RHS, it is a solution.
\textit{Graphing:} When $c = 0$, $y = \frac{4}{1-0} = 4$. This is a horizontal line at $y=4$. Looking at the provided graphs, Graph C correctly shows the horizontal line at $y=4$ for $c=0$, and vertical asymptotes (where $ce^x = 1 \implies x = -\ln(c)$) matching the behavior for $c > 0$.
\textbf{Correct Choice:} C

\subsection*{Problem 8}
\textbf{Verify solution and graph:} $\phi(x) = Ce^{4x} + 3$ for $\frac{dy}{dx} - 4y = -12$.

\textbf{Solution:}
First, find the derivative:
\begin{align*}
    \frac{dy}{dx} &= 4Ce^{4x}
\end{align*}
Substitute into the left side of the differential equation:
\begin{align*}
    \frac{dy}{dx} - 4y &= 4Ce^{4x} - 4(Ce^{4x} + 3) \\
    &= 4Ce^{4x} - 4Ce^{4x} - 12 \\
    &= -12
\end{align*}
This matches the right side, confirming it is a solution. 
\textit{Graphing:} When $C = 0$, $y = 3$. This forms a horizontal line. For $C > 0$, it is exponential growth, and for $C < 0$, it is exponential decay, but shifted up by 3. Graph C accurately models these characteristics.
\textbf{Correct Choice:} C

\subsection*{Problem 9}
\textbf{Find value of parameter \textit{m} for solution $\phi(x) = x^m$:}
If $\phi(x) = x^m$, then $\frac{dy}{dx} = mx^{m-1}$ and $\frac{d^2y}{dx^2} = m(m-1)x^{m-2}$.

\textbf{Part (a):} $5x^2\frac{d^2y}{dx^2} + 14x\frac{dy}{dx} - 2y = 0$
\begin{align*}
    5x^2(m(m-1)x^{m-2}) + 14x(mx^{m-1}) - 2(x^m) &= 0 \\
    x^m [5m(m-1) + 14m - 2] &= 0
\end{align*}
Assuming $x \neq 0$, we solve the characteristic polynomial:
\begin{align*}
    5m^2 - 5m + 14m - 2 &= 0 \\
    5m^2 + 9m - 2 &= 0 \\
    (5m - 1)(m + 2) &= 0 \implies m = \frac{1}{5}, -2
\end{align*}

\textbf{Part (b):} $x^2\frac{d^2y}{dx^2} + 3x\frac{dy}{dx} - 6y = 0$
\begin{align*}
    x^2(m(m-1)x^{m-2}) + 3x(mx^{m-1}) - 6(x^m) &= 0 \\
    x^m [m^2 - m + 3m - 6] &= 0 \\
    m^2 + 2m - 6 &= 0
\end{align*}
Use the quadratic formula:
\begin{align*}
    m &= \frac{-2 \pm \sqrt{4 - 4(1)(-6)}}{2} = \frac{-2 \pm \sqrt{28}}{2} \\
    &= \frac{-2 \pm 2\sqrt{7}}{2} = -1 \pm \sqrt{7}
\end{align*}
Answers: $-1 + \sqrt{7}, \quad -1 - \sqrt{7}$.

\subsection*{Problem 10}
\textbf{Consider $9y\frac{dy}{dx} - 16x = 0$ with implicit solutions $16x^2 - 9y^2 = C$. Initial condition $y(x_0)=0$.}

\textbf{Part (a): Does E\&U Theorem imply unique solution?}
Rewrite as $\frac{dy}{dx} = \frac{16x}{9y}$. Since $f(x,y) = \frac{16x}{9y}$ requires division by $y$, it is discontinuous at any point where $y=0$. Since our initial condition is $y(x_0)=0$, we are centered precisely on this discontinuity. The theorem does \textit{not} imply a unique solution.
\textbf{Correct Choice:} C

\textbf{Part (b): Show no solution in neighborhood of $x_0 \neq 0$.}
Substituting the condition $y=0$ at $x=x_0$ into the original DE gives:
\begin{align*}
    9(0)\frac{dy}{dx} - 16x_0 &= 0 \\
    -16x_0 &= 0 \implies x_0 = 0
\end{align*}
This directly contradicts the given assumption that $x_0 \neq 0$.

\textbf{Part (c): Show two distinct solutions for $y(0)=0$.}
Substitute the point $(0,0)$ into the implicit solution $16x^2 - 9y^2 = C$:
\[ 16(0)^2 - 9(0)^2 = C \implies C = 0 \]
Now substitute $C = 0$ back in and solve for $y$:
\begin{align*}
    16x^2 - 9y^2 &= 0 \\
    9y^2 &= 16x^2 \\
    y^2 &= \frac{16}{9}x^2 \\
    y &= \pm \frac{4}{3}x
\end{align*}
The two distinct explicit solutions are $y = \frac{4}{3}x$ and $y = -\frac{4}{3}x$.

\newpage
\section{Part 3: In-Depth Analysis of Problems and Techniques}

\subsection{A) Problem Types and General Approach}
The homework assignment focuses heavily on solution verification and theoretical guarantees. We can group the problems into three major categories:

\textbf{Type 1: Solution Verification (Problems 1, 2, 3, 4, 7, 8)}
\begin{itemize}
    \item \textit{Goal:} Prove that a given function (explicit or implicit) satisfies a given differential equation.
    \item \textit{Strategy (Explicit):} Compute the necessary derivatives of the proposed solution $\phi(x)$ and substitute both $\phi(x)$ and its derivatives directly into the DE. Simplify both sides algebraically to show $LHS = RHS$.
    \item \textit{Strategy (Implicit):} Use implicit differentiation on the given relationship to find $\frac{dy}{dx}$. If necessary, substitute the original relationship back in or use algebraic manipulations (like multiplying by $e^{-6xy}$ in Problem 4) to make your derivative match the given DE exactly.
\end{itemize}

\textbf{Type 2: Applying the Existence and Uniqueness Theorem (Problems 5, 6, 10a)}
\begin{itemize}
    \item \textit{Goal:} Determine if an initial value problem is mathematically guaranteed to have exactly one solution without actually solving it.
    \item \textit{Strategy:} Isolate $\frac{dy}{dx}$ to define $f(x,y)$. Compute $\frac{\partial f}{\partial y}$. Evaluate the continuity of both $f(x,y)$ and $\frac{\partial f}{\partial y}$ specifically at the initial coordinate $(x_0, y_0)$. If either has division by zero (like $y=0$), asymptotes, or square roots of negative numbers at that exact point, the theorem fails.
\end{itemize}

\textbf{Type 3: Method of Undetermined Parameters (Problem 9)}
\begin{itemize}
    \item \textit{Goal:} Find specific values of a constant $m$ such that a generic test function (like $y = x^m$) becomes a true solution to a DE.
    \item \textit{Strategy:} Compute $y'$ and $y''$ in terms of $m$ and substitute them into the DE. Factor out the common $x$ term, leaving an algebraic "characteristic equation" strictly in terms of $m$. Solve for $m$ using factoring or the quadratic formula.
\end{itemize}

\subsection{B) Key Algebraic and Calculus Manipulations}
\begin{itemize}
    \item \textbf{Implicit Differentiation with the Product Rule (Problem 1b):} When differentiating $xy^7$, you must treat it as $u = x, v = y^7$, yielding $y^7 \frac{dx}{dx} + x \cdot 7y^6 \frac{dy}{dx}$. A common mistake is to write this simply as $7xy^6 \frac{dy}{dx}$.
    \item \textbf{Fraction Decomposition / Common Denominators (Problem 7):} To verify the RHS, we expanded $\left(\frac{4}{1-ce^x}\right)^2$ and subtracted a related fraction. The crucial algebraic trick was multiplying the second term by $\frac{1-ce^x}{1-ce^x}$ to achieve a common denominator, which neatly collapsed the numerator to match the LHS.
    \item \textbf{Chain Rule with Negative Exponents (Problem 1, 7):} Converting roots to fractional exponents $\sqrt[8]{5-x} = (5-x)^{1/8}$ before differentiating prevents messy quotient rules. Remembering the inner derivative (multiplying by $-1$ because of the $-x$) is a paramount step.
    \item \textbf{Clever Algebraic Forcing (Problem 4):} After finding $\frac{dy}{dx} = \frac{x^2 - ye^{6xy}}{xe^{6xy} + y^2}$, we needed to match it to a formula containing $e^{-6xy}$. The "trick" was multiplying the fraction by a disguised form of 1: $\frac{e^{-6xy}}{e^{-6xy}}$, which perfectly cancelled the positive exponents and created the negative exponents required by the prompt.
\end{itemize}

\newpage
\section{Part 4: "Cheatsheet" and Tips for Success}

\subsection*{Summary of Important Formulas}
\textbf{The Existence and Uniqueness (E\&U) Test for $y' = f(x,y), y(x_0)=y_0$:}
\begin{enumerate}
    \item Identify $f(x,y)$.
    \item Compute $\frac{\partial f}{\partial y}$.
    \item Check both for continuity at $(x_0, y_0)$. (Look out for denominators that equal zero, or even-powered roots of negative numbers).
\end{enumerate}

\textbf{Euler's Characteristic Equation trick for $y = x^m$:}
For equations of the form $ax^2y'' + bxy' + cy = 0$ (Cauchy-Euler form), assuming $y=x^m$ leads to the characteristic quadratic equation: 
\[ am(m-1) + bm + c = 0 \]

\subsection*{Tips, Tricks, and Shortcuts}
\begin{itemize}
    \item \textbf{Shortcut for Partial Derivatives:} When calculating $\frac{\partial f}{\partial y}$, mentally treat $x$ as if it were a number like 5. So, $\frac{\partial}{\partial y}(x^2 y^3 + \sin(x)) = x^2(3y^2) + 0 = 3x^2y^2$.
    \item \textbf{Graphing Families of Solutions:} To determine which graph matches a solution like $y = Ce^{4x} + 3$, always test the trivial case $C=0$ first. This instantly eliminated options that didn't have a horizontal line at $y=3$. Then, test large limits ($x \to \infty$) to check for asymptotes.
    \item \textbf{Matching Implicit Derivatives:} If your calculated $\frac{dy}{dx}$ looks different than the textbook's answer, see if you can substitute the \textit{original implicit equation} back into your derivative to simplify it.
\end{itemize}

\subsection*{Common Pitfalls to Watch Out For}
\begin{itemize}
    \item \textbf{Forgetting the Inner Derivative (Chain Rule):} In $y=(5-x)^{1/8}$, many students write $\frac{1}{8}(5-x)^{-7/8}$ and stop. You must multiply by the derivative of $(5-x)$, which is $-1$.
    \item \textbf{Over-relying on the E\&U Theorem:} Remember, if $f$ or $\frac{\partial f}{\partial y}$ are discontinuous, the theorem says a unique solution \textit{is not guaranteed}. It does \textit{not} definitively say "no solution exists." There might be multiple solutions, or none. (e.g., Problem 10 showed multiple solutions existed at a discontinuity).
    \item \textbf{Ignoring Domain Restrictions:} A solution to a DE is only valid on an interval where both the function and its derivative are continuous and defined. Always check where denominators equal zero (as seen in Problem 1).
\end{itemize}

\newpage
\section{Part 5: Conceptual Synthesis and The "Big Picture"}

\subsection{A) Thematic Connections}
The core theme of this introductory chapter is \textbf{"Forward vs. Backward Thinking."} In standard Calculus (Calc I/II), you are given a function and asked to find its properties (area under the curve, rate of change, tangent line). Differential equations reverse this entirely. You are given the \textit{properties} (rates of change) and must find the \textit{function}. 
This mirrors the relationship between differentiation and integration. In fact, the simplest DEs, such as $\frac{dy}{dx} = f(x)$, are solved purely by integrating both sides: $y = \int f(x) dx$. Thus, solving a DE is fundamentally a generalized, more complex form of integration.

\subsection{B) Forward and Backward Links}
\textbf{Backward Links:} This topic relies completely on algebraic manipulation and differentiation rules learned in Calculus I, and the implicit differentiation techniques from Multivariable Calculus (Calculus III). Without mastery of the chain rule and product rule, solution verification is impossible.
\\
\textbf{Forward Links:} Learning to verify solutions and check Initial Value Problems is the bedrock for the rest of the course. Soon, you will transition from \textit{verifying} solutions to actually \textit{generating} them using techniques like Separation of Variables, Integrating Factors, and Laplace Transforms. Furthermore, the Existence and Uniqueness Theorem is your first introduction to qualitative analysis—understanding the behavior of a system without explicitly solving it, which is crucial for advanced nonlinear dynamics and chaos theory.

\newpage
\section{Part 6: Real-World Application and Modeling}

\subsection{A) Concrete Scenarios in Finance and Economics}
\textbf{1. Continuous Compound Interest (Basic Modeling):}
In finance, if an investment grows continuously, the rate of change of the account balance is directly proportional to the current balance. This is the simplest formulation of a differential equation in finance.

\textbf{2. Derivative Pricing and Geometric Brownian Motion:}
Advanced quantitative finance relies heavily on differential equations to model asset prices over time. Stock prices are often modeled using a Stochastic Differential Equation (SDE) known as Geometric Brownian Motion. The model balances a predictable drift (average return) with a stochastic (random) volatility factor. Understanding initial value problems is required to solve for the expected value path of the asset.

\textbf{3. The Black-Scholes Model:}
The Black-Scholes equation is a partial differential equation (PDE) that dictates the price of options contracts over time. Finding a unique solution to the Black-Scholes PDE requires specific "boundary conditions" and "initial conditions"—such as the payoff of the option at its expiration date.

\subsection{B) Model Problem Setup}
Let us set up a concrete mathematical model for a Retirement Drawdown Fund using ordinary differential equations. 
\begin{itemize}
    \item \textbf{Scenario:} An investor has a retirement account that generates a continuous fixed return, but they are also withdrawing money at a continuous fixed rate for living expenses. We want to find a function $V(t)$ that models the total value of the account over time.
    \item \textbf{Variables:} 
        \begin{itemize}
            \item $t$: Time in years.
            \item $V(t)$: The value of the retirement portfolio at time $t$.
            \item $r$: The annual continuous growth rate (e.g., $0.05$ for $5\%$).
            \item $W$: The continuous annual withdrawal rate (e.g., $\$40,000$ per year).
        \end{itemize}
    \item \textbf{The Initial Value Problem:}
    The rate of change of the portfolio $\frac{dV}{dt}$ is equal to the money gained from interest minus the money removed via withdrawals.
    \[ \frac{dV}{dt} = rV(t) - W \]
    The initial condition is the amount of money in the account at retirement ($t=0$):
    \[ V(0) = V_0 \quad \text{(e.g., } \$1,000,000 \text{)} \]
    This constitutes a first-order linear initial value problem that can be solved to determine exactly when (if ever) the investor will run out of money.
\end{itemize}

\newpage
\section{Part 7: Common Variations and Untested Concepts}

The homework set provided an excellent overview of verifying solutions and applying the Existence Theorem. However, standard curriculums for Section 1.2 typically include a few additional problem types:

\subsection{Missing Concept 1: Solving for the Constant \textit{C} in an IVP}
Your homework only asked you to verify general solutions or determine uniqueness. It did not explicitly test the mechanical process of finding the specific constant $C$ given a general solution and an initial condition (except briefly in 10c).

\textbf{Example Problem:} Given the general solution $y = Ce^{-3x} + x^2 - 1$ to a differential equation, find the particular solution that satisfies the initial condition $y(0) = 4$.

\textbf{Solution Strategy:} Substitute $x = 0$ and $y = 4$ into the general solution and solve for $C$.
\begin{align*}
    4 &= C e^{-3(0)} + (0)^2 - 1 \\
    4 &= C(1) - 1 \\
    C &= 5
\end{align*}
The particular solution is $y = 5e^{-3x} + x^2 - 1$.

\subsection{Missing Concept 2: Singular Solutions}
A singular solution is a solution to a differential equation that cannot be obtained from the general one-parameter family of solutions $y = f(x, C)$ by choosing any real number for $C$. 

\textbf{Example Problem:} The DE $\frac{dy}{dx} = xy^{1/2}$ has a general family of solutions $y = \left( \frac{x^2}{4} + C \right)^2$. However, by looking at the DE, we can see that if $y = 0$, then $\frac{dy}{dx} = 0$, making $y=0$ a perfectly valid solution ($0 = x \cdot 0$). 
Can we pick a constant $C$ to make the general solution equal $y=0$?
\[ \left( \frac{x^2}{4} + C \right)^2 = 0 \implies \frac{x^2}{4} + C = 0 \implies C = -\frac{x^2}{4} \]
Because $C$ must be a constant, it cannot equal $-\frac{x^2}{4}$. Thus, $y=0$ is a \textit{singular solution} that is "lost" when generating the standard family of curves.

\newpage
\section{Part 8: Advanced Diagnostic Testing: "Find the Flaw"}

Below are five problems with fully worked out, yet fundamentally flawed, solutions. Your task is to act as the grader, identify the specific mathematical error, explain why it is wrong in one sentence, and provide the correct steps.

\hrule
\vspace{0.3cm}
\textbf{Flawed Problem 1:} Verify that the implicit equation $x^2 + y^2 = 25$ is a solution to the differential equation $y' = -\frac{x}{y}$.
\\
\textbf{Flawed Solution:}
Differentiate both sides with respect to $x$:
\begin{align*}
    \frac{d}{dx}[x^2 + y^2] &= \frac{d}{dx}[25] \\
    2x + 2y \cdot y' &= 25 \\
    2y \cdot y' &= 25 - 2x \\
    y' &= \frac{25 - 2x}{2y}
\end{align*}
Since this does not match $-\frac{x}{y}$, it is not a solution.

\vspace{0.2cm}
\textit{Your Job: Find the flaw in the derivative of the constant.}
\vspace{0.3cm}
\hrule

\vspace{0.3cm}
\textbf{Flawed Problem 2:} Use the Existence and Uniqueness theorem to determine if the IVP $\frac{dy}{dx} = \sqrt{y}$ with $y(0) = 0$ has a unique solution.
\\
\textbf{Flawed Solution:}
Here, $f(x,y) = \sqrt{y}$. The initial point is $(0,0)$. Because the square root function $f(x,y) = \sqrt{y}$ is continuous at $y=0$, the function is continuous in a region around $(0,0)$. Therefore, the theorem guarantees a unique solution.

\vspace{0.2cm}
\textit{Your Job: Find the flaw regarding the partial derivative requirement.}
\vspace{0.3cm}
\hrule

\vspace{0.3cm}
\textbf{Flawed Problem 3:} Determine if $y = e^{2x}$ is a solution to $y'' - 4y = 0$.
\\
\textbf{Flawed Solution:}
Find the first and second derivatives:
\begin{align*}
    y &= e^{2x} \\
    y' &= 2e^{2x} \\
    y'' &= 2e^{2x} \quad \text{(applying chain rule again)}
\end{align*}
Substitute into the DE:
\[ 2e^{2x} - 4(e^{2x}) = -2e^{2x} \neq 0 \]
It is not a solution.

\vspace{0.2cm}
\textit{Your Job: Find the flaw in the second derivative calculation.}
\vspace{0.3cm}
\hrule

\vspace{0.3cm}
\textbf{Flawed Problem 4:} Determine the value of $m$ such that $y = e^{mx}$ is a solution to $y'' - 5y' + 6y = 0$.
\\
\textbf{Flawed Solution:}
Let $y = e^{mx}$. Then $y' = m e^{mx}$ and $y'' = m^2 e^{mx}$. 
Substitute into the DE:
\begin{align*}
    m^2 e^{mx} - 5m e^{mx} + 6(e^{mx}) &= 0 \\
    e^{mx}(m^2 - 5m) &= 0
\end{align*}
Dividing out $e^{mx}$ leaves $m^2 - 5m = 0$, which factors to $m(m-5)=0$. Thus $m=0, 5$.

\vspace{0.2cm}
\textit{Your Job: Find the dropped constant term during factoring.}
\vspace{0.3cm}
\hrule

\vspace{0.3cm}
\textbf{Flawed Problem 5:} Determine the interval of definition for the solution $y = \ln|x|$ to the DE $y' = \frac{1}{x}$, with initial condition $y(1) = 0$.
\\
\textbf{Flawed Solution:}
The function $y = \ln|x|$ is continuous and defined for all real numbers except $x=0$. The derivative $y' = \frac{1}{x}$ is also defined for all real numbers except $x=0$. Therefore, the interval of validity is $(-\infty, 0) \cup (0, \infty)$.

\vspace{0.2cm}
\textit{Your Job: Find the flaw regarding the definition of an interval for an IVP (it must be a single, unbroken interval containing the initial point).}
\vspace{0.3cm}
\hrule

\end{document}